\documentclass[aspectratio=169]{beamer}

% ---------- Beamer Theme ----------
\usetheme{metropolis} % or remove if you prefer your default
\setbeamertemplate{navigation symbols}{}
\setbeamertemplate{footline}[frame number]

% ---------- Packages ----------
\usepackage{amsmath, amssymb, bm, physics, mathtools}
\usepackage{siunitx}
\usepackage{tikz}
\usepackage{booktabs}
\usepackage{graphicx}

% ---------- Macros ----------
\newcommand{\uu}{\bm{u}}
\newcommand{\n}{\bm{n}}
\newcommand{\x}{\bm{x}}
\newcommand{\I}{\mathbf{I}}
\newcommand{\D}{\mathsf{D}} % rate-of-strain tensor
% \DeclareMathOperator{\grad}{\nabla}
\DeclareMathOperator{\divg}{\nabla\!\cdot}
%\DeclareMathOperator{\curl}{\nabla\times}

\title{From Euler to Navier--Stokes}
\subtitle{Material derivative, forces, and the viscous stress tensor}
\author{Laura A. Miller}
\institute{University of Arizona}
\date{}

\begin{document}

% =========================================================
\begin{frame}
  \titlepage
\end{frame}

% =========================================================

\begin{frame}{Gentle Introductions to Fluid Dynamics}
\textbf{For students new to fluids, including biologists:}
\tiny
\begin{itemize}
  \item \textbf{MIT OpenCourseWare: Introduction to Fluid Mechanics (1.060)}  
  Conceptual, with physical intuition before equations  
  \url{https://ocw.mit.edu/courses/1-060-engineering-mechanics-ii-fluid-dynamics-fall-2013/}

  \item \textbf{Harvard Natural Sciences Lectures: Fluids \& Biological Flows}  
  Demonstrations + biological context  
  \url{https://projects.iq.harvard.edu/sciencelectures}

  \item \textbf{NPTEL “Fluid Mechanics for Biological Systems” (IIT Guwahati)}  
  Directly aimed at life-science audiences  
  \url{https://nptel.ac.in/courses/102103084}
\end{itemize}

\end{frame}

\begin{frame}{More Advanced Fluid Dynamics Lecture Resources — University Courses}
\begin{itemize}
  \item \textbf{MIT OpenCourseWare}
  \begin{itemize}
    \item 2.25 Advanced Fluid Mechanics  
    \url{https://ocw.mit.edu/courses/2-25-advanced-fluid-mechanics-fall-2013/}
    \item 2.06 Fluid Dynamics  
    \url{https://ocw.mit.edu/courses/2-06-fluid-dynamics-spring-2020/}
  \end{itemize}

  \item \textbf{Cambridge University Engineering Dept}  
  Fluid Mechanics I \& II lecture notes  
  \url{https://www-civ.eng.cam.ac.uk/cued-lecture-notes/fluids/}

  \item \textbf{Caltech GALCIT — Fluid Mechanics}  
  \url{https://www.galcit.caltech.edu/education}
\end{itemize}

\end{frame}

\begin{frame}{Video Lecture Playlists}

\begin{itemize}
  \item \textbf{Stanford CFD Lectures (Duraisamy)}  
  \url{https://www.youtube.com/playlist?list=PLd3hlSJsX_ImKP68wfKjvRlKq2Ot3PuuX}

  \item \textbf{NPTEL – IIT Kanpur Fluid Mechanics}  
  \url{https://nptel.ac.in/courses/112104118}

  \item \textbf{Flow Visualization — Univ. of Colorado Boulder}  
  \url{https://www.flowvis.org/}

\end{itemize}

\end{frame}

\begin{frame}{Additional References}

\begin{itemize}
  \item \textbf{L.C. Evans — PDE Notes (for NS theory)}  
  \url{https://math.ucdavis.edu/~evans/}

  \item \textbf{Dabiri Lab — Bio-inspired Fluid Dynamics}  
  \url{https://dabirilab.com/videos/}

  \item \textbf{APS Division of Fluid Dynamics Tutorials}  
  \url{https://www.aps.org/units/dfd/meetings/video.cfm}
\end{itemize}

\end{frame}


% --- INTRODUCTORY SLIDES: Notation and Vector Operators ---
\section{Summary of Notation}
\footnotesize
\begin{frame}{Basic Definitions}
\begin{block}{Position vector}
\[
\mathbf{x} = (x_1, x_2, x_3) = (x, y, z)
\]
\end{block}
\begin{block}{Velocity vector}
\[
\mathbf{u} = (u_1, u_2, u_3) = (u, v, w)
\]
\end{block}
\begin{block}{Other quantities}
\begin{itemize}
  \item $\rho$ – density (kg/m$^3$)
  \item $p$ – pressure (Pa)
  \item $\mu$ – dynamic viscosity (Pa·s)
  \item $\mathbf{g}$ – gravitational acceleration vector (m/s$^2$)
\end{itemize}
\end{block}
\end{frame}

%------------------------------------------------------------

\begin{frame}{Vector differential operators}
\begin{columns}[T,onlytextwidth]
\column{0.55\textwidth}
\begin{block}{Gradient of a scalar $\phi$}
\[
\nabla \phi =
\left(
\frac{\partial \phi}{\partial x},
\frac{\partial \phi}{\partial y},
\frac{\partial \phi}{\partial z}
\right)
\]
\end{block}

\begin{block}{Divergence of a vector $\mathbf{u}$}
\[
\nabla \!\cdot\! \mathbf{u}
= \frac{\partial u}{\partial x}
+ \frac{\partial v}{\partial y}
+ \frac{\partial w}{\partial z}
\]
\end{block}

\column{0.45\textwidth}
\begin{block}{Curl of a vector $\mathbf{u}$}
\[
\nabla \!\times\! \mathbf{u} =
\begin{vmatrix}
\mathbf{i} & \mathbf{j} & \mathbf{k} \\
\dfrac{\partial}{\partial x} &
\dfrac{\partial}{\partial y} &
\dfrac{\partial}{\partial z} \\
u & v & w
\end{vmatrix}
=
\left[
\begin{aligned}
\frac{\partial w}{\partial y} - \frac{\partial v}{\partial z} \\
\frac{\partial u}{\partial z} - \frac{\partial w}{\partial x} \\
\frac{\partial v}{\partial x} - \frac{\partial u}{\partial y}
\end{aligned}
\right]
\]
\end{block}
\end{columns}
\end{frame}


%------------------------------------------------------------

\begin{frame}{Summary of notation}
\begin{center}
\renewcommand{\arraystretch}{1.4}
\begin{tabular}{l l}
\textbf{Symbol} & \textbf{Meaning} \\ \hline
$\mathbf{x} = (x, y, z)$ & Position vector \\
$\mathbf{u} = (u, v, w)$ & Velocity vector \\
$p$ & Pressure \\
$\rho$ & Density \\
$\mu$ & Dynamic viscosity \\
$\mathbf{g}$ & Body-force acceleration \\
$\nabla \phi$ & Gradient of scalar $\phi$ \\
$\nabla \!\cdot\! \mathbf{u}$ & Divergence of $\mathbf{u}$ \\
$\nabla \!\times\! \mathbf{u}$ & Curl (vorticity) of $\mathbf{u}$ \\
$\nabla^2 \mathbf{u}$ & Laplacian (componentwise) of $\mathbf{u}$
\end{tabular}
\end{center}
\end{frame}

\section{Kinematics \& Material Derivative}

\begin{frame}{Derivation of the Material Derivative}
\small
\textbf{Goal:} Express the rate of change of a field $\phi(\bm{x},t)$ \\
as experienced by a moving fluid parcel.

\vspace{0.4cm}

\textbf{1. Follow a fluid parcel:}
\[
\bm{x} = \bm{x}(t), \qquad \frac{d\bm{x}}{dt} = \bm{u}(\bm{x},t)
\]

\textbf{2. Apply the chain rule along the trajectory:}
\[
\frac{d\phi}{dt}
= \pdv{\phi}{t} + \frac{d\bm{x}}{dt} \cdot \grad \phi
\]

\textbf{3. Substitute the parcel velocity:}
\[
\boxed{
\frac{D\phi}{Dt}
= \pdv{\phi}{t} + (\bm{u}\cdot\grad)\phi
}
\]

\vspace{0.4cm}
\end{frame}

\begin{frame}{Derivation of the Material Derivative Cont'd}
\textbf{Interpretation:}
\begin{itemize}
  \item $\displaystyle \pdv{\phi}{t}$ — local (Eulerian) rate of change at a point.
  \item $\displaystyle (\bm{u}\cdot\grad)\phi$ — convective rate of change due to motion.
  \item Sum gives the total rate of change following the fluid.
\end{itemize}

\end{frame}


\begin{frame}{Material Derivative for a Scalar Field}

\textbf{Goal:} Rate of change of a quantity $\phi(x,y,z,t)$ following the flow.

\vspace{0.4cm}

A parcel moves as $\bm{x}(t) = (x(t), y(t), z(t))$ with velocity
\[
\frac{d\bm{x}}{dt} = (u, v, w) = \bm{u}.
\]

Apply the chain rule along the parcel trajectory:
\[
\frac{d\phi}{dt}
= \pdv{\phi}{t}
+ \frac{dx}{dt}\pdv{\phi}{x}
+ \frac{dy}{dt}\pdv{\phi}{y}
+ \frac{dz}{dt}\pdv{\phi}{z}.
\]

\vspace{0.4cm}

Substitute $\frac{dx}{dt}=u$, $\frac{dy}{dt}=v$, $\frac{dz}{dt}=w$:
\[
\boxed{
\frac{D\phi}{Dt}
= \pdv{\phi}{t}
+ u\pdv{\phi}{x}
+ v\pdv{\phi}{y}
+ w\pdv{\phi}{z}
}
\]

\end{frame}

\begin{frame}{Material Derivative for a Vector Field}

For a vector field $\bm{u}(\bm{x},t) = (u,v,w)$,
apply the material derivative to each component:

\[
\frac{D\bm{u}}{Dt}
= \pdv{\bm{u}}{t} + (\bm{u}\cdot\nabla)\bm{u}.
\]

\vspace{0.4cm}

Expanded form:
\[
(\bm{u}\cdot\nabla)\bm{u}
= \left(
u\pdv{u}{x} + v\pdv{u}{y} + w\pdv{u}{z},\;
u\pdv{v}{x} + v\pdv{v}{y} + w\pdv{v}{z},\;
u\pdv{w}{x} + v\pdv{w}{y} + w\pdv{w}{z}
\right).
\]

\vspace{0.4cm}

\textbf{Interpretation:}
\begin{itemize}
  \item $\displaystyle \pdv{\bm{u}}{t}$ — local acceleration.
  \item $\displaystyle (\bm{u}\cdot\nabla)\bm{u}$ — convective acceleration.
\end{itemize}

\end{frame}

\begin{frame}{Component Form of the Material Derivative}

\textbf{Each component separately:}
\begin{align*}
\frac{Du}{Dt} &= \pdv{u}{t} + u\pdv{u}{x} + v\pdv{u}{y} + w\pdv{u}{z},\\[6pt]
\frac{Dv}{Dt} &= \pdv{v}{t} + u\pdv{v}{x} + v\pdv{v}{y} + w\pdv{v}{z},\\[6pt]
\frac{Dw}{Dt} &= \pdv{w}{t} + u\pdv{w}{x} + v\pdv{w}{y} + w\pdv{w}{z}.
\end{align*}

\vspace{0.3cm}

\textbf{Summary:}
\[
\boxed{
\frac{D\bm{u}}{Dt}
= \left(
\frac{Du}{Dt},
\frac{Dv}{Dt},
\frac{Dw}{Dt}
\right)
}
\]

\vspace{0.3cm}

These terms represent the total acceleration of a fluid parcel.

\end{frame}

\begin{frame}{Material Derivative}
We track a parcel with position $\x(t)$ and velocity $\uu(\x,t)$.
\[
\frac{D}{Dt} = \pdv{}{t} + \uu\cdot\grad
\]
Applied to velocity:
\[
\frac{D\uu}{Dt} = \pdv{\uu}{t} + (\uu\cdot\grad)\uu
\]
This is the acceleration of a fluid parcel (left-hand side of Euler / Navier--Stokes).
\end{frame}

\begin{frame}{Physical Meaning of Incompressibility}

\textbf{Imagine:} a small dyed blob of fluid moves with the flow.

\vspace{0.4cm}

As the blob moves, its shape may distort, but for an
\textbf{incompressible fluid} the \textit{volume of the blob stays constant}.

\vspace{0.4cm}

\textbf{Let:}
\begin{itemize}
  \item $V(t)$ — volume of the blob at time $t$
  \item $\bm{u}$ — velocity field of the fluid
\end{itemize}

Then incompressibility means:
\[
\boxed{\frac{dV}{dt} = 0.}
\]

\vspace{0.4cm}

We will now express this volume change in terms of the velocity field.
\end{frame}


\begin{frame}{Volume Flux Through a Closed Surface}
\small
Consider a fixed closed surface $S$ that encloses a region $V$ in the fluid,
with outward unit normal $\bm{n}$.

\vspace{0.4cm}

At each point on $S$, the component of velocity along the outward normal is
\[
\bm{u}\cdot\bm{n}.
\]

The volume of fluid leaving through a small surface element $\delta S$
in unit time is:
\[
(\bm{u}\cdot\bm{n})\,\delta S.
\]

\vspace{0.4cm}

Integrating over the whole surface gives the \textbf{net volume rate}
at which fluid leaves the region:
\[
\boxed{\frac{dV}{dt} = \oint_S \bm{u}\cdot\bm{n}\, dS.}
\]

If the fluid is incompressible, this rate must be zero for every closed surface.

\[
\boxed{\oint_S \bm{u}\cdot\bm{n}\, dS = 0.}
\]

\end{frame}

\begin{frame}{From Surface Flux to Divergence Condition}

\textbf{Recall the Divergence Theorem:}
\[
\boxed{
\oint_S \bm{u}\cdot\bm{n}\, dS
= \int_V (\nabla\cdot\bm{u})\, dV.
}
\]


If $\displaystyle \oint_S \bm{u}\cdot\bm{n}\, dS = 0$ for any closed surface $S$,
then the integrand must vanish everywhere inside $V$:

\[
\boxed{\nabla\cdot\bm{u} = 0.}
\]

\textbf{Interpretation:}
\begin{itemize}
  \item No net flux of volume out of any region.
  \item Fluid elements neither expand nor contract.
  \item The flow field is divergence-free.
\end{itemize}

\end{frame}


\begin{frame}{Mass Conservation}
For density $\rho(\x,t)$:
\[
\pdv{\rho}{t} + \divg(\rho \uu) = 0.
\]
Incompressible limit (constant $\rho$):
\[
\divg \uu = 0.
\]
\end{frame}



% =========================================================
\section{Forces: Pressure via Surface Tractions}

\begin{frame}{Pressure via Surface Traction}

\textbf{Setup:} Let $S$ be a closed surface enclosing a small dyed blob of fluid,
with outward unit normal $\bm{n}$ (pointing \emph{out of the blob}).

\vspace{0.35cm}

\textbf{Pressure assumption (isotropic normal stress):}
\[
\text{Force on } \delta S \text{ from the surrounding fluid } = \; p(\bm{x},t)\,\bm{n}\,\delta S,
\]
where $p(\bm{x},t)$ is a scalar, independent of the direction of $\bm{n}$ (\emph{the pressure}).

\vspace{0.35cm}
\end{frame}

\begin{frame}{Pressure via Surface Traction}
\textbf{Sign convention / traction on the blob:}
Because pressure is \emph{compressive}, the traction \emph{on the blob} is
\[
\boxed{\;\bm{t}_{\text{on blob}} = -\,p\,\bm{n}\;}
\]
(the outward normal $\bm{n}$ points out of the blob, so the pressure force points inward).

\vspace{0.35cm}

\textbf{Net pressure force on the blob:}
\[
\boxed{\;\bm{F}_{p,\text{blob}} = \displaystyle -\oint_S p\,\bm{n}\,dS\;}
\]
(The negative sign encodes that compression acts opposite the outward normal.)
\end{frame}

\begin{frame}{From Surface Traction to $-\nabla p$ (Volume Form)}

\textbf{Gradient/Divergence step (a.k.a. “gradient theorem” here):}
Apply Gauss’s divergence theorem \emph{componentwise}.
For each Cartesian basis vector $\bm{e}_i$,
\[
\oint_S (p\,\bm{e}_i)\cdot\bm{n}\,dS
= \int_V \nabla\cdot(p\,\bm{e}_i)\,dV
= \int_V \partial_i p\,dV.
\]
Stacking components gives the vector identity
\[
\boxed{\;\oint_S p\,\bm{n}\,dS \;=\; \int_V \nabla p\,dV\;}
\]
(“surface integral of $p\,\bm{n}$ equals volume integral of $\nabla p$”).

\vspace{0.35cm}

\textbf{Hence, the net pressure force on the blob:}
\[
\bm{F}_{p,\text{blob}} = -\oint_S p\,\bm{n}\,dS
= -\int_V \nabla p\,dV.
\]

\end{frame}

\begin{frame}{From Surface Traction to $-\nabla p$ (Volume Form)}
\textbf{Small-blob limit (nearly uniform $\nabla p$):}
If the blob is small so that $\nabla p$ is approximately constant in $V$,
\[
\boxed{\;\bm{F}_{p,\text{blob}} \approx -\,(\nabla p)\,\delta V\;}
\]
points toward lower pressure.

\vspace{0.35cm}

\textbf{Sign check (Newton’s III):}
The force the \emph{blob exerts on the surrounding fluid} is the opposite:
\[
\bm{F}_{p,\text{surroundings}} \approx +\,(\nabla p)\,\delta V.
\]
\end{frame}


% =========================================================
\section{Euler’s Equations (Inviscid)}

\begin{frame}{Equations of Motion for an Ideal Fluid}

\textbf{Definition (ideal fluid):}
\begin{enumerate}
  \item[(i)] \textit{Incompressible}: no dyed blob of fluid can change volume as it moves.
  \item[(ii)] \textit{Constant density}: $\rho$ is the same for all fluid elements and for all time $t$.
  \item[(iii)] \textit{Pressure traction}: the force exerted across a geometrical surface element
               $\bm{n}\,\delta S$ within the fluid is
               \[
                 p\,\bm{n}\,\delta S,
               \]
               where $p(x,y,z,t)$ is a scalar, independent of the normal direction $\bm{n}$ (the \emph{pressure}).
\end{enumerate}

\medskip
Forces are compressive, acting opposite the outward normal on a dyed blob.

\end{frame}

\begin{frame}{Momentum Balance on a Dyed Blob with Gravity}

\textbf{Blob:} small control mass of volume $\delta V$ moving with the fluid; outward normal $\bm{n}$.

\medskip
\textbf{Pressure force on the blob:}
\[
\bm{F}_p \;=\; -\oint_S p\,\bm{n}\,dS \;=\; -\int_{\delta V} \nabla p\,dV \;\approx\; -(\nabla p)\,\delta V.
\]

\textbf{Body force (gravity) per unit mass:} $\bm{g}$ \quad $\Rightarrow$ \quad
\[
\bm{F}_g \;=\; \int_{\delta V} \rho\,\bm{g}\,dV \;\approx\; \rho\,\bm{g}\,\delta V.
\]
\end{frame}

\begin{frame}{Momentum Balance on a Dyed Blob with Gravity}
\textbf{Total force on the blob:}
\[
\boxed{\;\bm{F}_{\text{total}} \;=\; \big(-\nabla p + \rho\,\bm{g}\big)\,\delta V\;}
\]

\textbf{Newton’s 2nd law (blob mass $\rho\,\delta V$):}
\[
\rho\,\delta V\,\frac{D\bm{u}}{Dt} \;=\; \big(-\nabla p + \rho\,\bm{g}\big)\,\delta V.
\]

Cancel $\delta V$ to obtain the acceleration equation.

\end{frame}

\begin{frame}{Euler’s Equations for an Ideal (Inviscid) Incompressible Fluid}

From the blob momentum balance:
\[
\boxed{\;\frac{D\bm{u}}{Dt} \;=\; -\,\frac{1}{\rho}\,\nabla p \;+\; \bm{g}\;}
\quad\text{with}\quad
\frac{D}{Dt}=\pdv{}{t}+\bm{u}\cdot\nabla.
\]

\medskip
\textbf{Incompressibility (kinematic):}
\[
\boxed{\;\nabla\cdot\bm{u}=0\;}
\]
(Fluid parcels do not change volume; density is constant.)

\end{frame}

\begin{frame}{Euler (Inviscid) Equations in Components}

Let $\bm{u}=(u,v,w)$ and $\bm{g}=(g_x,g_y,g_z)$.
Then
\begin{align*}
\frac{Du}{Dt} &= \pdv{u}{t} + u\pdv{u}{x} + v\pdv{u}{y} + w\pdv{u}{z}
               \;=\; -\,\frac{1}{\rho}\,\pdv{p}{x} + g_x,\\[6pt]
\frac{Dv}{Dt} &= \pdv{v}{t} + u\pdv{v}{x} + v\pdv{v}{y} + w\pdv{v}{z}
               \;=\; -\,\frac{1}{\rho}\,\pdv{p}{y} + g_y,\\[6pt]
\frac{Dw}{Dt} &= \pdv{w}{t} + u\pdv{w}{x} + v\pdv{w}{y} + w\pdv{w}{z}
               \;=\; -\,\frac{1}{\rho}\,\pdv{p}{z} + g_z.
\end{align*}

\medskip
\textbf{Incompressibility:}
\[
\boxed{\;\pdv{u}{x}+\pdv{v}{y}+\pdv{w}{z}=0\;}
\]

\end{frame}


% =========================================================
\section{From Euler to Navier--Stokes: Add Viscosity}

\begin{frame}{Momentum Balance with Stresses}
Cauchy momentum equation introduces the Cauchy stress tensor $\bm{\sigma}$:
\[
\rho \frac{D\uu}{Dt} = \divg \bm{\sigma} + \rho \bm{f}.
\]
Split stress into pressure and deviatoric parts:
\[
\bm{\sigma} = -p \I + \bm{\tau},
\]
where $\bm{\tau}$ is the viscous (deviatoric) stress tensor.
\end{frame}


\begin{frame}{The Stress Vector}
\small
Let $\bm{x}$ denote the position of a fixed point in the fluid,
and let $\delta S$ be a small geometrical surface element through $\bm{x}$,
with unit normal $\bm{n}$.

\vspace{0.4cm}

The \textbf{force} exerted by the fluid on the side of the surface
towards which $\bm{n}$ is directed is assumed to be
\[
\boxed{\;\bm{t}\,\delta S\;}
\]
where $\bm{t}$ is the \textbf{stress vector} acting on that surface.

\vspace{0.4cm}

For an inviscid fluid:
\[
\bm{t} = -\,p(\bm{x},t)\,\bm{n}.
\]

In general, $\bm{t}$ may have both normal and tangential components.

\vspace{0.3cm}
\textbf{Goal:} Express $\bm{t}$ in terms of quantities defined at the point $\bm{x}$,
independent of the orientation of the surface.
\end{frame}
\begin{frame}{Definition of the Stress Tensor}
\small
To describe all possible stress vectors at a point,
we define nine local quantities $T_{ij}$:

\vspace{0.3cm}

\[
\boxed{\;
T_{ij} = \text{$i$-component of stress on a surface with normal in the $j$-direction.}
\;}
\]

\vspace{0.3cm}

In matrix form (Cartesian coordinates):
\[
\bm{T} =
\begin{pmatrix}
T_{xx} & T_{xy} & T_{xz}\\
T_{yx} & T_{yy} & T_{yz}\\
T_{zx} & T_{zy} & T_{zz}
\end{pmatrix}.
\]

Each column corresponds to a face normal direction;
each row corresponds to a component of force.

\vspace{0.4cm}

Our goal is to relate the stress vector $\bm{t}$ on an arbitrarily oriented
surface (with unit normal $\bm{n}$) to these tensor components $T_{ij}$.
\end{frame}
\begin{frame}{Stress on an Arbitrary Surface}
\footnotesize
Consider a small tetrahedron of fluid (Fig.~6.3 in Acheson):
one large face has outward normal $\bm{n}$,
and three coordinate faces are perpendicular to $x$, $y$, and $z$ axes.

\vspace{0.3cm}

Apply the principle of linear momentum to the tetrahedron:

\begin{itemize}
  \item Stress on main face (normal $\bm{n}$): $\bm{t}\,\delta S$.
  \item Stress on coordinate faces (normals $-\bm{e}_x$, $-\bm{e}_y$, $-\bm{e}_z$): 
  $-\,T_{i1}n_x\delta S$, $-\,T_{i2}n_y\delta S$, $-\,T_{i3}n_z\delta S$.
\end{itemize}

Balancing $i$-components of force as the tetrahedron shrinks ($\delta V\propto L^3$, $\delta S\propto L^2$):

\[
t_i = T_{ij}\,n_j,
\quad \text{(summation on $j$ implied).}
\]

\vspace{0.4cm}

\textbf{Vector form:}
\[
\boxed{\;\bm{t} = \bm{T}\,\bm{n}\;}
\]

This relation is known as \textbf{Cauchy’s stress principle}.

\end{frame}
\begin{frame}{Cauchy’s Equation of Motion}
\small
Consider any infinitesimal fluid element of volume $\delta V$.

\textbf{Forces acting:}
\begin{itemize}
  \item Surface (stress) forces: $\displaystyle \oint_S \bm{t}\,dS = \oint_S \bm{T}\bm{n}\,dS$
  \item Body forces: $\rho\,\bm{g}\,\delta V$
\end{itemize}

\vspace{0.3cm}

Apply the divergence theorem:
\[
\oint_S \bm{T}\bm{n}\,dS = \int_V \nabla\cdot\bm{T}\,dV.
\]

\textbf{Linear momentum balance:}
\[
\rho\,\frac{D\bm{u}}{Dt} = \nabla\cdot\bm{T} + \rho\,\bm{g}.
\]

\vspace{0.4cm}

\textbf{This is Cauchy’s equation of motion.}

\vspace{0.2cm}

Special cases:
\[
\bm{T} = -p\mathbf{I} \quad \Rightarrow \quad
\rho\frac{D\bm{u}}{Dt} = -\nabla p + \rho\bm{g} \quad (\text{Euler}).
\]

\end{frame}


% =========================================================
\section{How to Build the Viscous Stress Tensor}

\begin{frame}{Cauchy Stress and Momentum Balance}
\begin{block}{Traction vector on an arbitrary surface}
For a surface with unit normal $\mathbf{n}$ at a point,
\[
\mathbf{t}(\mathbf{n}) = \boldsymbol{\sigma}\,\mathbf{n},
\qquad
\boldsymbol{\sigma} = (\sigma_{ij})\ \text{is the Cauchy stress tensor.}
\]
Here $\boldsymbol{\sigma}$ maps the direction $\mathbf{n}$ to the
force per unit area across that surface.
\end{block}

\begin{block}{Cauchy equation of motion}
Balance of linear momentum for a fluid of density $\rho$:
\[
\rho\,\frac{D\mathbf{u}}{Dt}
=
\nabla\!\cdot \boldsymbol{\sigma}
+ \rho\,\mathbf{g}.
\]
\end{block}
\end{frame}

\begin{frame}{Split Stress into Pressure and Viscous Parts}
\begin{columns}[T,onlytextwidth]
\column{0.55\textwidth}
\[
\boldsymbol{\sigma} = -p\,I + \boldsymbol{\tau},
\]
\begin{itemize}
  \item $-pI$ is the isotropic (pressure) stress.
  \item $\boldsymbol{\tau}$ is the deviatoric (viscous) stress.
\end{itemize}

\column{0.45\textwidth}
\begin{block}{Symmetry}
No internal torques (no couple stresses) implies
\[
\sigma_{ij}=\sigma_{ji}.
\]
\end{block}
\end{columns}
\end{frame}

\begin{frame}{The Rate-of-Strain Tensor: What It Measures}
\small
\textbf{Goal:} Describe how nearby fluid particles deform relative to each other.

\medskip

For a velocity field $\mathbf{u}(\mathbf{x})$, the velocity gradient $\nabla \mathbf{u}$
contains \emph{all} local information about stretching, shearing, and rotation.

\medskip

But only the \textbf{symmetric part} produces physical deformation:
\[
D = \frac12\!\left( \nabla \mathbf{u} + (\nabla \mathbf{u})^{\!\top} \right).
\]

\begin{itemize}
\item $D$ measures how a small fluid element is \emph{stretched or sheared}.
\item The antisymmetric part
\[
W = \frac12\!\left( \nabla\mathbf{u} - (\nabla\mathbf{u})^{\!\top} \right)
\]
represents rigid-body rotation only — it does \emph{not} generate stresses.
\item Thus $D$ is the unique second-order tensor describing deformation rate.
\end{itemize}

\end{frame}

\begin{frame}{Geometric Interpretation of $D$}

Consider two nearby particles separated by a small vector $d\mathbf{x}$.

\[
\text{relative velocity} \;\;\; d\mathbf{u} = (\nabla \mathbf{u})\, d\mathbf{x}.
\]

Decompose using symmetry:
\[
d\mathbf{u}
= D\,d\mathbf{x} + W\,d\mathbf{x}.
\]

\begin{itemize}
\item $D\,d\mathbf{x}$: changes length and angle of $d\mathbf{x}$  
      (stretching and shearing).
\item $W\,d\mathbf{x}$: rotates the line element without deformation.
\end{itemize}

\vspace{0.2cm}

Thus $D$ captures the \emph{actual deformation} of the material, making it
the natural quantity for determining viscous stresses.

\end{frame}

\begin{frame}{What Constitutive Laws Should Look Like}

A constitutive law for viscous stress $\boldsymbol{\tau}$ must satisfy:

\begin{enumerate}
\item \textbf{Linearity} (Newtonian assumption):  
      stress depends linearly on the rate of deformation.
\item \textbf{Locality}: stress at a point depends only on $\nabla \mathbf{u}$ at that point.
\item \textbf{Objectivity}: rigid-body motions must not generate stress.
\item \textbf{Isotropy}: material has no preferred direction.
\end{enumerate}

\medskip

These restrictions dramatically narrow down the possible forms of $\boldsymbol{\tau}$.

\end{frame}

\begin{frame}{Why Isotropy Determines the Functional Form}

For an \textbf{isotropic} material, $\boldsymbol{\tau}$ cannot depend on direction.
Therefore, $\boldsymbol{\tau}$ may be built only from:

\begin{itemize}
\item the symmetric tensor $D$,
\item its trace, $\nabla\!\cdot\mathbf{u}$ (a scalar),
\item and the identity tensor $I$.
\end{itemize}

\medskip

\textbf{Most general linear, isotropic tensor function:}
\[
\boxed{
\boldsymbol{\tau}
=
2\mu\, D
+
\lambda\,(\nabla\!\cdot\mathbf{u})\,I
}
\]

\begin{itemize}
\item $\mu$ — dynamic viscosity (shear response)
\item $\lambda$ — bulk viscosity (volumetric response)
\end{itemize}

\textbf{This is the only possible linear isotropic relation.}

\end{frame}

\begin{frame}{Special Case: Incompressible Newtonian Fluids}

If the flow is incompressible:
\[
\nabla\!\cdot\mathbf{u}=0,
\]
the constitutive law simplifies to
\[
\boldsymbol{\tau}=2\mu D.
\]

Component form:
\[
\sigma_{ij}
=
-p\,\delta_{ij}
+
\mu\!\left(
\frac{\partial u_i}{\partial x_j}
+
\frac{\partial u_j}{\partial x_i}
\right).
\]

\begin{itemize}
\item Pressure accounts for isotropic normal stresses.
\item Viscosity produces shear stresses proportional to deformation rate.
\end{itemize}

\end{frame}

\begin{frame}{Divergence of the Viscous Stress}

For constant $\mu$ (and $\lambda$ if used),
\[
\nabla\!\cdot(2\mu D)
=
\mu \nabla^2 \mathbf{u}
+
\mu\,\nabla(\nabla\!\cdot\mathbf{u}).
\]

\medskip

For incompressible flows:
\[
\nabla\!\cdot\mathbf{u}=0
\quad\Rightarrow\quad
\nabla\!\cdot\boldsymbol{\tau}
=
\mu \nabla^2 \mathbf{u}.
\]

This provides the familiar \textbf{viscous diffusion term} in Navier–Stokes.

\end{frame}

\begin{frame}{Putting It All Together: Navier–Stokes}

Start from Cauchy momentum balance:
\[
\rho\,\frac{D\mathbf{u}}{Dt}
=
\nabla\!\cdot \boldsymbol{\sigma}
+
\rho\,\mathbf{g},
\qquad
\boldsymbol{\sigma}
=
-pI + \boldsymbol{\tau}.
\]

Insert the constitutive law for incompressible Newtonian fluids:
\[
\boldsymbol{\tau}=2\mu D,
\qquad
\nabla\!\cdot\boldsymbol{\tau}
=\mu\nabla^2\mathbf{u}.
\]

\[
\boxed{
\rho\left(
\frac{\partial\mathbf{u}}{\partial t}
+
(\mathbf{u}\!\cdot\!\nabla)\mathbf{u}
\right)
=
-\nabla p
+
\mu\nabla^2\mathbf{u}
+
\rho\mathbf{g},
\qquad
\nabla\!\cdot\mathbf{u}=0.
}
\]

\textbf{This is the Navier–Stokes equation for an incompressible Newtonian fluid.}

\end{frame}









% =========================================================
\section{Navier--Stokes Equations}

\begin{frame}{Navier--Stokes Equations in Partial Derivative Form}
\footnotesize
For an incompressible Newtonian fluid with constant viscosity $\mu$ and density $\rho$,
the momentum equations in $(x,y,z)$ coordinates are:

\[
\rho \left(
\frac{\partial u}{\partial t}
+ u\,\frac{\partial u}{\partial x}
+ v\,\frac{\partial u}{\partial y}
+ w\,\frac{\partial u}{\partial z}
\right)
=
-\frac{\partial p}{\partial x}
+ \mu\left(
\frac{\partial^2 u}{\partial x^2}
+ \frac{\partial^2 u}{\partial y^2}
+ \frac{\partial^2 u}{\partial z^2}
\right)
+ \rho g_x,
\]

\[
\rho \left(
\frac{\partial v}{\partial t}
+ u\,\frac{\partial v}{\partial x}
+ v\,\frac{\partial v}{\partial y}
+ w\,\frac{\partial v}{\partial z}
\right)
=
-\frac{\partial p}{\partial y}
+ \mu\left(
\frac{\partial^2 v}{\partial x^2}
+ \frac{\partial^2 v}{\partial y^2}
+ \frac{\partial^2 v}{\partial z^2}
\right)
+ \rho g_y,
\]

\[
\rho \left(
\frac{\partial w}{\partial t}
+ u\,\frac{\partial w}{\partial x}
+ v\,\frac{\partial w}{\partial y}
+ w\,\frac{\partial w}{\partial z}
\right)
=
-\frac{\partial p}{\partial z}
+ \mu\left(
\frac{\partial^2 w}{\partial x^2}
+ \frac{\partial^2 w}{\partial y^2}
+ \frac{\partial^2 w}{\partial z^2}
\right)
+ \rho g_z.
\]

\medskip

Incompressibility constraint:
\[
\frac{\partial u}{\partial x}
+ \frac{\partial v}{\partial y}
+ \frac{\partial w}{\partial z}
= 0.
\]

\end{frame}


\begin{frame}{Navier--Stokes in Component Form}
\footnotesize
For an incompressible Newtonian fluid with constant viscosity $\mu$ and density $\rho$,
the momentum equations in component form are:

\[
\rho\left(
\frac{\partial u_i}{\partial t}
+ u_j\,\frac{\partial u_i}{\partial x_j}
\right)
=
-\frac{\partial p}{\partial x_i}
+ \mu\,\frac{\partial^2 u_i}{\partial x_j^2}
+ \rho\,g_i,
\qquad i = 1,2,3.
\]

\medskip

Explicitly, in $(x,y,z)$ coordinates:

\[
\begin{aligned}
\rho\left(u_{1,t} + u_1 u_{1,x} + u_2 u_{1,y} + u_3 u_{1,z}\right)
&= -p_{,x} + \mu\left(u_{1,xx}+u_{1,yy}+u_{1,zz}\right) + \rho g_1,\\[4pt]
\rho\left(u_{2,t} + u_1 u_{2,x} + u_2 u_{2,y} + u_3 u_{2,z}\right)
&= -p_{,y} + \mu\left(u_{2,xx}+u_{2,yy}+u_{2,zz}\right) + \rho g_2,\\[4pt]
\rho\left(u_{3,t} + u_1 u_{3,x} + u_2 u_{3,y} + u_3 u_{3,z}\right)
&= -p_{,z} + \mu\left(u_{3,xx}+u_{3,yy}+u_{3,zz}\right) + \rho g_3.
\end{aligned}
\]

The incompressibility condition is

\[
u_{1,x} + u_{2,y} + u_{3,z} = 0.
\]

\end{frame}

\begin{frame}{Understanding Each Term: A Guided Breakdown}

The momentum equation in component form:

\[
\rho\left(
\frac{\partial u_i}{\partial t}
+ u_j\,\frac{\partial u_i}{\partial x_j}
\right)
=
-\frac{\partial p}{\partial x_i}
+ \mu\,\frac{\partial^2 u_i}{\partial x_j^2}
+ \rho\,g_i.
\]

\begin{itemize}
\item $\displaystyle \frac{\partial u_i}{\partial t}$ — \textbf{local acceleration}
\item $\displaystyle u_j \frac{\partial u_i}{\partial x_j}$ — \textbf{advective acceleration}
\item $\displaystyle -\frac{\partial p}{\partial x_i}$ — \textbf{pressure-gradient force}
\item $\displaystyle \mu \frac{\partial^2 u_i}{\partial x_j^2}$ — \textbf{viscous diffusion of momentum}
\item $\rho g_i$ — \textbf{body forces} (gravity, buoyancy, etc.)
\end{itemize}

\medskip

This is just Newton’s $F=ma$ applied to a fluid element, written in tensor form.

\end{frame}

\begin{frame}{Steady Navier--Stokes Equations}

For a steady incompressible Newtonian flow, 
\[
\frac{\partial \mathbf{u}}{\partial t} = 0.
\]

The Navier--Stokes equations reduce to:
\[
\rho\,(\mathbf{u}\cdot\nabla)\mathbf{u}
=
-\nabla p
+ \mu\,\nabla^2 \mathbf{u}
+ \rho\,\mathbf{g},
\qquad
\nabla\cdot\mathbf{u} = 0.
\]

\begin{itemize}
  \item No time dependence — acceleration comes only from advection.
  \item Important for fully developed pipe/channel flows.
  \item Still nonlinear and generally difficult to solve.
\end{itemize}

\end{frame}

\begin{frame}{1D Navier--Stokes (Incompressible)}

Assume velocity only in one direction: 
\[
\mathbf{u} = (u(x,t),\,0,\,0),
\qquad
u_y = u_z = 0.
\]

Then incompressibility gives:
\[
\frac{\partial u}{\partial x} = 0
\quad\Rightarrow\quad
u = u(t) \text{ or } u = u(y) \text{ in channel flow.}
\]

Momentum equation becomes:
\[
\rho\,\frac{\partial u}{\partial t}
=
-\frac{\partial p}{\partial x}
+ \mu\,\frac{\partial^2 u}{\partial y^2}
+ \rho\,g_x.
\]

\begin{itemize}
  \item Forms basis of Poiseuille and Couette flow.
  \item Nonlinearity disappears — no advection term.
  \item Widely used for analytical benchmark solutions.
\end{itemize}

\end{frame}

\begin{frame}{2D Navier--Stokes (Incompressible)}

Assume $\mathbf{u} = (u(x,y,t),\, v(x,y,t),\, 0)$.

Momentum equations:
\[
\rho\left(
\frac{\partial u}{\partial t}
+ u\,u_x + v\,u_y
\right)
=
- p_x
+ \mu\left(u_{xx}+u_{yy}\right)
+ \rho g_x,
\]

\[
\rho\left(
\frac{\partial v}{\partial t}
+ u\,v_x + v\,v_y
\right)
=
- p_y
+ \mu\left(v_{xx}+v_{yy}\right)
+ \rho g_y.
\]

Incompressibility constraint:
\[
u_x + v_y = 0.
\]

\begin{itemize}
  \item Used for cavity flow, cylinder wake, airfoil sections.
  \item Still nonlinear — supports vortex shedding and separation.
\end{itemize}

\end{frame}


% =========================================================
\section{Boundary Conditions}

\begin{frame}{Boundary Conditions for Navier--Stokes}

To solve the Navier--Stokes equations, the PDEs must be supplemented with
\emph{boundary conditions} on the domain boundary $\partial\Omega$.

\medskip

Boundary conditions determine:
\begin{itemize}
  \item how the fluid interacts with solid walls,
  \item what enters and exits the domain,
  \item whether momentum or stress is prescribed,
  \item whether the flow is constrained or free to adjust.
\end{itemize}

Different choices produce very different physical behavior.

\end{frame}

\begin{frame}{Dirichlet Boundary Conditions}

A \textbf{Dirichlet boundary condition} specifies the value of the solution:

\[
\mathbf{u} = \mathbf{u}_b \quad \text{on } \partial\Omega_D.
\]

Common examples:
\begin{itemize}
  \item \textbf{No-slip wall:}
        \[
        \mathbf{u} = \mathbf{0}
        \]
  \item \textbf{Prescribed inflow:}
        \[
        \mathbf{u} = \mathbf{u}_{\text{in}}(t)
        \]
  \item \textbf{Moving boundary:}
        \[
        \mathbf{u} = \mathbf{U}_{\text{wall}}(t)
        \]
\end{itemize}

Dirichlet BCs fix the velocity and completely constrain motion at the boundary.

\end{frame}

\begin{frame}{Neumann Boundary Conditions}
\footnotesize

A \textbf{Neumann boundary condition} specifies the \emph{flux} or \emph{gradient}
rather than the value of the field.

For momentum, this means prescribing the \textbf{traction}:

\[
\boldsymbol{\sigma}\,\mathbf{n}
=
\mathbf{t}_b
\quad \text{on } \partial\Omega_N.
\]

Examples:
\begin{itemize}
  \item \textbf{Stress-free surface:}
        \[
        \boldsymbol{\sigma}\,\mathbf{n} = \mathbf{0}
        \]
  \item \textbf{Outflow traction:}
        \[
        -p\,\mathbf{n} + \mu (\nabla\mathbf{u} + \nabla\mathbf{u}^{\!\top})\mathbf{n}
        = \mathbf{t}_{\text{out}}
        \]
  \item \textbf{No-penetration (normal only):}
        \[
        \mathbf{u}\cdot\mathbf{n} = 0
        \quad\Rightarrow\quad
        \text{normal velocity gradient unspecified}
        \]
\end{itemize}

Neumann BCs arise naturally from momentum balance and variational formulations.

\end{frame}

\begin{frame}{Robin (Mixed) Boundary Conditions}
\footnotesize
A \textbf{Robin boundary condition} combines Dirichlet and Neumann data:

\[
\alpha\,\mathbf{u}
+
\beta\,(\boldsymbol{\sigma}\,\mathbf{n})
=
\mathbf{g}
\quad \text{on } \partial\Omega_R.
\]

Examples:
\begin{itemize}
  \item \textbf{Navier slip} (partial slip at walls):
        \[
        \mathbf{u}_\parallel
        =
        \ell_s \left(
        (\boldsymbol{\sigma}\,\mathbf{n})_\parallel
        \right)
        \]
        where $\ell_s$ is a slip length.
  \item \textbf{Porous or resistive boundaries} (Darcy-type):
        \[
        \mathbf{u}\cdot\mathbf{n}
        +
        K\,(\boldsymbol{\sigma}\,\mathbf{n})\cdot\mathbf{n}
        =
        0
        \]
\end{itemize}

Robin BCs interpolate between fully constrained and fully free boundaries.


\end{frame}

\begin{frame}{Other Common Boundary Conditions in Fluid Mechanics}

\begin{itemize}
  \item \textbf{Slip (inviscid wall assumption):}
        \[
        \mathbf{u}\cdot\mathbf{n} = 0,
        \qquad
        (\boldsymbol{\sigma}\,\mathbf{n})_\parallel = \mathbf{0}
        \]
  \item \textbf{Periodic boundaries:}
        \[
        \mathbf{u}(\mathbf{x}) = \mathbf{u}(\mathbf{x} + L_i \mathbf{e}_i),
        \qquad
        p(\mathbf{x}) = p(\mathbf{x} + L_i \mathbf{e}_i)
        \]
        used in channel flow and turbulence.
  \item \textbf{Far-field / open boundaries:}
        \[
        \frac{\partial \mathbf{u}}{\partial n} = 0
        \quad\text{or}\quad
        \boldsymbol{\sigma}\,\mathbf{n} = \mathbf{0}
        \]
  \item \textbf{Free surface (air–water interface):}
        \[
        \boldsymbol{\sigma}\,\mathbf{n}
        = -p_{\text{atm}}\mathbf{n}
        + \gamma\,\kappa\,\mathbf{n}
        \quad (\text{surface tension})
        \]
\end{itemize}

Different boundary choices can drastically change the flow solution.

\end{frame}

\begin{frame}{How to Choose Boundary Conditions}

Consider:
\begin{itemize}
  \item Physics of the domain (solid wall? open flow? interface?)
  \item What is known and what the solution must determine
  \item Numerical stability (outflow BCs are delicate)
  \item Whether the boundary should constrain momentum or stress
\end{itemize}

\medskip

\textbf{Rule of thumb:}
\begin{itemize}
  \item Dirichlet for inflow and solid walls,
  \item Neumann/traction or ``do–nothing'' for outflow,
  \item Periodic when the physics repeats.
\end{itemize}

\end{frame}

\begin{frame}{How Many Boundary Conditions Do We Need?}

For incompressible Navier--Stokes, the velocity satisfies
second--order PDEs. Therefore:

\begin{itemize}
  \item Each velocity component needs \emph{one} boundary condition
        on the boundary $\partial\Omega$.
  \item At every boundary point, specify \textbf{either}
        \begin{itemize}
           \item a \textbf{Dirichlet} condition (velocity), e.g.\ $\mathbf{u}=\mathbf{u}_b$, \\
           \item or a \textbf{Neumann} condition (traction), e.g.\ $\boldsymbol{\sigma}\mathbf{n}=\mathbf{t}_b$.
        \end{itemize}
  \item Velocity BCs may be mixed over the boundary:
        \[
        \partial\Omega
        = \partial\Omega_D \cup \partial\Omega_N,
        \qquad
        \partial\Omega_D \cap \partial\Omega_N = \varnothing.
        \]
  \item The pressure has no evolution equation and is determined only up to a constant:
        \begin{itemize}
           \item one additional constraint is required (e.g.\ fix $p$ at a point, or prescribe traction on part of the boundary).
        \end{itemize}
\end{itemize}

\end{frame}

\begin{frame}{Why We Do Not Prescribe Pressure Directly}
\footnotesize
In incompressible Navier--Stokes, the pressure has \textbf{no time evolution equation}.
Instead, it is determined by the momentum balance and the incompressibility constraint.

\medskip

\begin{itemize}
  \item The pressure appears only through its \textbf{gradient}:
  \[
    -\nabla p
  \]
  so adding a constant to $p$ does not change the physics.
  \item Pressure is therefore a \textbf{Lagrange multiplier} enforcing
  \[
    \nabla\!\cdot \mathbf{u} = 0.
  \]
  \item The pressure field is obtained by solving a Poisson equation:
  \[
    \nabla^2 p = \nabla\!\cdot\big( \rho(\mathbf{u}\!\cdot\!\nabla)\mathbf{u} \big),
  \]
  which uses boundary information from the \textbf{velocity} or \textbf{traction}, not $p$ itself.
  \item Only one extra condition is needed (e.g.\ fix $p$ at a point or specify traction on part of $\partial\Omega$) to make $p$ unique.
\end{itemize}

\end{frame}


\begin{frame}{Do People Ever Prescribe Pressure as a Boundary Condition?}
\footnotesize
\textbf{Yes — but only in specific situations.}  
In incompressible Navier--Stokes, pressure is not an independent field with its own
evolution equation, so we do \emph{not} freely assign $p$ on boundaries.

\medskip

\textbf{Cases often called ``pressure BCs'':}
\begin{itemize}
  \item \textbf{Traction (Neumann) condition:}
        \[
        \boldsymbol{\sigma}\mathbf{n} = -p_{\text{ext}}\mathbf{n}
        \]
        e.g.\ free surface, outflow to atmosphere.
  \item \textbf{Reference value for uniqueness:}
        \[
        p(x_0)=0
        \]
        fixes the additive constant.
  \item \textbf{Special cases:}
        steady Stokes flow, certain outflow boundaries,
        or compressible flow where $p$ is thermodynamic.
\end{itemize}

\medskip

\textbf{Important:}  
Prescribing both $p$ \emph{and} normal velocity overconstrains the system and is ill-posed.

\end{frame}

\begin{frame}{Well-Posed vs.\ Ill-Posed Boundary Conditions}

\textbf{Well-posed examples:}
\footnotesize
\begin{itemize}
  \item \textbf{Inflow–outflow:}
  \[
    \mathbf{u} = \mathbf{u}_{\text{in}} \text{ on } \partial\Omega_{\text{in}}, 
    \qquad
    \boldsymbol{\sigma}\mathbf{n} = \mathbf{0} \text{ on } \partial\Omega_{\text{out}}.
  \]
  \item \textbf{No-slip walls + traction elsewhere:}
  \[
    \mathbf{u} = \mathbf{0} \text{ on } \partial\Omega_{\text{wall}}, \quad
    \boldsymbol{\sigma}\mathbf{n} = \mathbf{t}_b \text{ on } \partial\Omega_{\text{free}}.
  \]
  \item \textbf{Periodic domain:}
  velocity and pressure equal on opposite faces.
\end{itemize}

\medskip

\textbf{Ill-posed or incorrect examples:}
\begin{itemize}
  \item Prescribing \emph{both} $\mathbf{u}$ and $\boldsymbol{\sigma}\mathbf{n}$ on the same boundary segment.
  \item Specifying pressure \emph{and} normal velocity at an outflow.
  \item Applying $\mathbf{u}=\mathbf{0}$ on the \emph{entire} boundary with no forcing
        $\Rightarrow$ trivial solution only.
\end{itemize}

\medskip

\textbf{Key rule:}  
Assign \emph{one} condition per velocity component at each boundary point—either Dirichlet or Neumann, not both.

\end{frame}



% =========================================================
\section{References}
\begin{frame}{References}
D.\,J. Acheson, \textit{Elementary Fluid Dynamics}. Clarendon Press.\\[4pt]
Any standard continuum mechanics text for the Cauchy stress and Newtonian constitutive law.
\end{frame}

% =========================================================
\section{Appendix: Mini-Derivations}

\begin{frame}{Mini: From Control Volume to Differential Form}
Apply Reynolds Transport Theorem and the divergence theorem to convert surface tractions to volume forces:
\[
\int_S \bm{\sigma}\,\n\,dS = \int_V \divg \bm{\sigma}\, dV.
\]
Combine with linear momentum balance to get the Cauchy equation.
\end{frame}

\begin{frame}{Mini: Identity for $\divg(2\mu \D)$}
For constant $\mu$:
\[
\divg(2\mu \D)
= \mu \nabla^2 \uu + \mu\,\grad(\divg \uu).
\]
(Use $\nabla^2 \uu \equiv \grad(\divg\uu) - \curl(\curl\uu)$ and symmetry of $\D$.)
\end{frame}

\end{document}
